\IfFormatAtLeastTF{2025-06-01}{\DocumentMetadata{lang=en-US,pdfversion=2.0,tagging=on}}{}
\documentclass[11pt]{amsart}
\usepackage[T1]{fontenc}
\IfFileExists{lmodern.sty}{\usepackage{lmodern}}{\PassOptionsToPackage{expansion=false}{microtype}}
\usepackage{amsmath,amssymb,amsthm,mathtools}
\usepackage{etoolbox}
\usepackage{tikz}
\usetikzlibrary{positioning}
\usepackage[nopatch=footnote]{microtype}
\usepackage{xcolor}
\usepackage[colorlinks=true,linkcolor=blue,citecolor=blue,urlcolor=blue]{hyperref}
\newif\ifanonymous
\anonymousfalse
\ifanonymous
  \newcommand{\paperauthor}{Anonymous}
  \newcommand{\coderepo}{the repository accompanying the submission}
\else
  % Repository locator for the public (non-anonymous) build of
% non_mf_groups_exist.tex.  Kept in a separate file so that the anonymous
% build has an anonymous source archive and not merely an anonymous PDF:
% exclude this file from a double-blind submission and the manuscript still
% compiles, printing "supplied with this submission as supplementary
% material" in place of the link.
%
% Loaded only when \anonymousfalse is set in the preamble.
\newcommand{\coderepository}{https://github.com/SauersML/group-approximation}
\newcommand{\coderepositoryname}{github.com/SauersML/group-approximation}

  \newcommand{\paperauthor}{\paperauthorname}
  \newcommand{\coderepo}{\url{https://github.com/SauersML/group-approximation}}
\fi
\hypersetup{pdftitle={[DRAFT] MF Recognition is Pi02-Complete},pdfauthor={\paperauthor},pdfsubject={Draft manuscript: arithmetical complexity of MF recognition},pdfkeywords={draft, MF groups, arithmetical hierarchy, finitely presented groups, HNN extensions, Higman embedding}}
\newcommand{\doi}[1]{\href{https://doi.org/#1}{\nolinkurl{doi:#1}}}
% Machine-readable source metadata attaching a printed statement to an
% accompanying Lean~4 declaration.  The macro is intentionally invisible in
% the PDF; the web edition and audit scripts consume its two arguments.
% Each \leanverified marker identifies a selected printed statement and a Lean
% declaration intended to prove it.  The markers do not assert line-by-line
% verification of the prose, sentence-census coverage, or formal verification
% of the headline recognition theorem.
\newcommand{\leanverified}[2]{\ignorespaces}
\allowdisplaybreaks
\raggedbottom
\emergencystretch=2em
\newtheorem{theorem}{Theorem}[section]
\newtheorem{proposition}[theorem]{Proposition}
\newtheorem{lemma}[theorem]{Lemma}
\newtheorem{corollary}[theorem]{Corollary}
\theoremstyle{remark}
\newtheorem{remark}[theorem]{Remark}
\theoremstyle{plain}
\theoremstyle{plain}
\newtheorem{mainthm}{Theorem}
\renewcommand{\themainthm}{\Alph{mainthm}}
\newcommand{\F}{\mathbb F}
\newcommand{\C}{\mathbb C}
\newcommand{\GL}{\operatorname{GL}}
\newcommand{\EL}{\operatorname{EL}}
\newcommand{\U}{\mathcal U}
\newcommand{\Rad}{\operatorname{Rad}}
\newcommand{\tr}{\operatorname{tr}}
\newcommand{\opnorm}[1]{\left\lVert #1\right\rVert}
\newcommand{\hsnorm}[1]{\left\lVert #1\right\rVert_2}
\newcommand{\normal}[1]{\langle\!\langle #1\rangle\!\rangle}

\makeatletter
\patchcmd{\abstract}{\item[\hskip\labelsep\scshape\abstractname.]}{\item[]}{}{\PackageError{self-compression}{Could not remove abstract heading}{}}
\makeatother
\newcommand{\draftnotice}{%
  \begin{center}
  \fcolorbox{red!70!black}{red!5}{%
    \parbox{0.88\linewidth}{\centering\bfseries\color{red!70!black}
      DRAFT --- NOT READY FOR PUBLIC CONSUMPTION\par
      \smallskip
      \normalfont\color{black}This manuscript is actively being revised and
      has not been released for public consumption.}}
  \end{center}
  \medskip}

\begin{document}
\frenchspacing
\title{MF Recognition is $\Pi^0_2$-Complete}
\author{\paperauthor}
\subjclass[2020]{Primary 03D40; Secondary 20F10, 46L05, 03D80}
\keywords{MF groups, arithmetical hierarchy, finitely presented groups, HNN extensions, Higman embedding, norm matrix coronas}

\begin{abstract}
A countable group is MF if it admits finite-dimensional unitary models
that are asymptotically multiplicative and asymptotically separate
nonidentity elements in operator norm.  We show that deciding whether a
finite presentation defines an MF group is $\Pi^0_2$-complete.  So it is
impossible to construct an algorithm that decides from a finite
presentation whether the group it defines is MF, and impossible even to
enumerate the presentations of MF groups, or those of non-MF groups.  A
presentation defines an MF group if and only if, for every accuracy,
finitely many unitary matrices satisfy its relators to within that
accuracy while keeping prescribed nontrivial words away from the
identity; whether such matrices exist is decidable, by Tarski's theorem.  We also
construct, computably from any program, a finite presentation that
defines an MF group if and only if the program halts on infinitely many
inputs.  If the program halts on only finitely many inputs, the presented
group contains the non-MF group of~\cite{NonMF}; if it halts on
infinitely many inputs, the group is MF by a permanence theorem for HNN
extensions whose edge isomorphism is implemented by a unitary of a norm
matrix corona.
\end{abstract}

\maketitle
\draftnotice

\section{Introduction}

Not every countable group is MF: \emph{Non-MF Groups}~\cite{NonMF}
constructs a finitely generated simple group $H=\EL_{12}(L_{\F_2}(1,2))$
with property~\textup{(T)} such that every homomorphism from $H$ to an MF
group is trivial.  We turn to the recognition problem for finitely presented groups.

\begin{mainthm}[MF recognition is $\Pi^0_2$-complete]\label{thm:recognition}
Deciding whether a finite presentation defines an MF group is
$\Pi^0_2$-complete, and the complementary problem is
$\Sigma^0_2$-complete.  In particular it is impossible to construct an
algorithm that decides from a finite presentation whether the group it
defines is MF, or even one that enumerates the presentations of MF
groups or those of non-MF groups.  There is a computable map
$e\mapsto\widehat R_e$ such that $\widehat R_e$ presents an MF group if
and only if the $e$-th partial computable function has infinite domain.
\end{mainthm}

\begin{figure}[htb]
\centering
\begin{tikzpicture}[font=\small,
  cls/.style={draw, rounded corners=2pt, inner sep=4pt},
  probs/.style={font=\small, align=center, inner sep=1pt}]
\node[cls] (d1) at (0,0)      {$\Delta^0_1$};
\node[cls] (s1) at (-2.0,1.5) {$\Sigma^0_1$};
\node[cls] (p1) at (2.0,1.5)  {$\Pi^0_1$};
\node[cls] (s2) at (-2.0,3.2) {$\Sigma^0_2$};
\node[cls] (p2) at (2.0,3.2)  {$\Pi^0_2$};
\draw[thin] (d1)--(s1);
\draw[thin] (d1)--(p1);
\draw[thin] (s1)--(s2);
\draw[thin] (p1)--(p2);
\draw[thin] (s1)--(p2);
\draw[thin] (p1)--(s2);
\node[probs, above=2mm of s2]
  {$\mathrm{FIN}$\\ $\boldsymbol{\mathrm{NONMF}_{\mathrm{fp}}}$};
\node[probs, above=2mm of p2]
  {$\mathrm{INF}$; $\mathrm{TOT}$\\ torsion-freeness\\
   $\boldsymbol{\mathrm{MF}_{\mathrm{fp}}}$};
\node[probs, anchor=east] at (-3.0,1.5) {halting\\ problem};
\node[probs, below=2mm of d1]
  {word problem in $H$ (Lemma~\ref{lem:seed})};
\end{tikzpicture}
\caption{The arithmetical hierarchy.  Lines indicate inclusions, and each
listed problem is complete at its level; the word problem in $H$ is
decidable.  The two sets from Theorem~\ref{thm:recognition} appear in
bold.  $\mathrm{MF}_{\mathrm{fp}}\in\Pi^0_2$ by finite certificates
(Section~\ref{sec:local}).  The
family $\widehat R_e$ reduces $\mathrm{INF}$ to
$\mathrm{MF}_{\mathrm{fp}}$ and $\mathrm{FIN}$ to
$\mathrm{NONMF}_{\mathrm{fp}}$ (Section~\ref{sec:compiler}).}
\label{fig:levels}
\end{figure}

For each $n$, a certificate consists of a dimension $d$, a label for every
word of length at most $n$, and a normal-closure proof for every word
labelled trivial.  It is valid if there are unitaries
$U_1,\dots,U_k\in\U(d)$ such that $\opnorm{r_i(U)-1}\le2^{-n}$ for each
relator $r_i$ and $\opnorm{w(U)-1}\ge1/4$ for every word labelled
nontrivial.  These conditions are polynomial equations and inequalities
over the real numbers, so validity is decidable by Tarski's theorem.  A
finite presentation defines an MF group if and only if it has
a valid certificate for every $n$.  So
$\mathrm{MF}_{\mathrm{fp}}\in\Pi^0_2$ and
$\mathrm{NONMF}_{\mathrm{fp}}\in\Sigma^0_2$
(Section~\ref{sec:local}).

Given a program index $e$, we construct a finite presentation
$\widehat R_e$ in Section~\ref{sec:compiler}.  If the $e$-th partial function
has finite domain, then the presented group contains $H$ and is not MF\@.
If its domain is infinite, the edge isomorphism in the defining HNN
extension is implemented by a unitary of a norm matrix corona, and the
presented group is MF by Theorem~\ref{thm:hnn-permanence}.
Korchagin proved that the MF property is closed under direct
limits~\cite[Proposition~13]{Korchagin}.  So every finitely
generated non-MF group is a quotient of a finitely presented one.

Proofs appear here at the level of their ideas; complete proofs of every
statement are formalized in the Lean proof assistant and checked by its
kernel.

\subsection*{Relation to prior work}

Dehn initiated the study of decision problems for finitely presented
groups.  Adian and Rabin later proved that every Markov property is
undecidable from a finite presentation~\cite{Adian,Rabin}.  A property is
Markov if it holds for some finitely presented group and fails for every
group containing some fixed finitely presented group.  The trivial group
is MF, whereas every group containing $H$ is non-MF.  Thus MF is a Markov
property, and the Adian--Rabin theorem proves undecidability.  Our theorem
determines the precise level of undecidability.

Lempp proved that recognizing torsion-free finitely presented groups is
$\Pi^0_2$-complete~\cite{Lempp}.  For recursively presented groups,
Bilanovic--Chubb--Roven proved that every Markov property is
$\Pi^0_2$-hard to recognize~\cite{BCR}; their construction begins our
reduction.  We must then replace a recursive presentation by a finite
one without changing whether the group is MF.  A general Higman embedding
does not preserve the MF property.  So we use HNN extensions
chosen to satisfy the permanence theorem proved below.

\section{MF groups}\label{sec:prelim}

A countable group $G$ is \emph{MF} if there are positive
integers $d_n$ and maps $V_n\colon G\to\U(d_n)$, with $V_n(1)=1$, such that
\[
 \opnorm{V_n(gh)-V_n(g)V_n(h)}\longrightarrow0
 \qquad(g,h\in G),
\]
and
\[
 \limsup_n\opnorm{V_n(g)-1}>0
 \qquad(g\in G\setminus\{1\}).
\]
This is the operator norm notion of MF group introduced by
Carri\'on--Dadarlat--Eckhardt~\cite{CDE}.  Here MF group
always means this operator-norm local-approximation property.  The
strong-convergence convention also requires the matrix models to recover
the norms of the left regular representation~\cite{GaoEtAl,Schafhauser}.
Equivalently, $G$ embeds in the
unitary group of the norm matrix corona associated to the sequence
$\mathbf d=(d_n)$,
\[
 \mathcal Q_{\mathbf d}
 =\prod_nM_{d_n}(\C)\big/\bigoplus_nM_{d_n}(\C),
\]
where $\bigoplus_nM_{d_n}(\C)$ is the $c_0$-direct sum of the sequences with
$\opnorm{x_n}\to0$.  A separable $C^*$-algebra is called MF if it embeds as
a $C^*$-subalgebra of a norm matrix corona~\cite{BK}.  Restricting models
to a subgroup shows that subgroups of MF groups are MF\@.
\leanverified{Sofic/MFDefinitions}{GroupApproximation.standardMFDefinitions_iff}
\leanverified{Sofic/OperatorMFPositiveControls}{GroupApproximation.IsOperatorMF.subgroup}

Throughout, the binary Leavitt algebra $L_{\F_2}(1,2)$ is the unital
$\F_2$-algebra generated by $s_0,s_1,t_0,t_1$ subject to $t_is_j=\delta_{ij}$
and $s_0t_0+s_1t_1=1$~\cite{Leavitt}, and
$H=\EL_{12}(L_{\F_2}(1,2))$ is the subgroup of
$\GL_{12}(L_{\F_2}(1,2))$ generated by the elementary matrices
$e_{ij}(a)=1+aE_{ij}$, $i\ne j$.  By~\cite[Theorem~B]{NonMF}, every homomorphism from
$H$ to an MF group is trivial; in particular $H$ is not MF\@.

\section{Locality and finite certificates}\label{sec:local}

We use the following finite-presentation form of Korchagin's local
criterion~\cite[Propositions~7--8]{Korchagin}.

\begin{lemma}[Korchagin's criterion]\label{lem:local-models}
Let $G=\langle x_1,\ldots,x_k\mid r_1,\ldots,r_m\rangle$.  Then $G$ is MF
if and only if the following holds.  Given $\varepsilon>0$ and finitely
many words $w_1,\ldots,w_q$ representing nonidentity elements of $G$,
there are $d\ge1$ and $U_1,\ldots,U_k\in\U(d)$ such that
\[
 \opnorm{r_i(U)-1}\le\varepsilon\quad(1\le i\le m),
 \qquad
 \opnorm{w_j(U)-1}\ge\tfrac12\quad(1\le j\le q).
\]
\leanverified{Sofic/OperatorMFLocalNormalization}{GroupApproximation.OperatorMFLocalNormalization.isOperatorMF_iff_isNormApproximable_one}
\end{lemma}

\begin{proof}
Korchagin proves the criterion with one chosen word representing each
element of a finite subset of $G$, and with separation constant
$1$~\cite[Propositions~7--8]{Korchagin}.  If $\omega_j$ is his chosen
representative of the element represented by $w_j$, then
$w_j\omega_j^{-1}$ is a finite product of conjugates of the relators and
their inverses.  So, for the finite family under consideration, a sufficiently
small relator tolerance preserves separation at least $1/2$ at the
prescribed words.  Conversely, Korchagin's
amplification argument replaces any fixed positive separation constant
by $1$ while changing the relator tolerance only by a constant factor.
\end{proof}

A set $A\subseteq\mathbb N$ belongs to $\Pi^0_2$ if there is a decidable
relation $R$ such that
\[
 x\in A\quad\Longleftrightarrow\quad
 \forall n\,\exists c\;R(x,n,c).
\]
Its complement then belongs to $\Sigma^0_2$; see~\cite{Soare}.  Fix a
computable coding of finite presentations by natural numbers.  From a
code $P$ we can recover the generators and relators of the presentation.
Write $G_P$ for the presented group and set
\[
 \mathrm{MF}_{\mathrm{fp}}=\{P:G_P\text{ is MF}\},
 \qquad
 \mathrm{NONMF}_{\mathrm{fp}}=\mathbb N\setminus
 \mathrm{MF}_{\mathrm{fp}}.
\]

\begin{proposition}[finite certificates]\label{prop:mf-upper-bound}
There is a decidable predicate $C(P,n,c)$ such that
\[
 P\in\mathrm{MF}_{\mathrm{fp}}
 \quad\Longleftrightarrow\quad
 \forall n\,\exists c\;C(P,n,c).
\]
So
$\mathrm{MF}_{\mathrm{fp}}\in\Pi^0_2$ and
$\mathrm{NONMF}_{\mathrm{fp}}\in\Sigma^0_2$.
\end{proposition}

\begin{proof}
Fix
$P=\langle x_1,\ldots,x_k\mid r_1,\ldots,r_m\rangle$, and let $W_n$ be
the finite set of reduced words of length at most $n$.  A certificate
contains the following data:
\begin{enumerate}
\item a positive integer $d$;
\item a partition $W_n=T\sqcup S$;
\item for each $w\in T$, an expression of $w$ as a product of conjugates
of the relators and their inverses.
\end{enumerate}
The certificate is valid if every expression in item~3 freely reduces to
the indicated word and there are $U_1,\ldots,U_k\in\U(d)$ satisfying
\[
 \opnorm{r_i(U)-1}\le2^{-n}\quad(1\le i\le m),
 \qquad
 \opnorm{w(U)-1}\ge\tfrac14\quad(w\in S).
\]
Here $w(U)$ denotes evaluation at the matrices $U_j$, with
$x_j^{-1}$ evaluated as $U_j^*$.

Free reduction decides whether the expressions in item~3 are correct.
So every word admitted to $T$ is trivial in $G_P$.  It remains to
decide whether the required matrices exist.  Write $U_j=A_j+iB_j$ with
real matrices $A_j$ and $B_j$.  Unitarity becomes a finite system of
polynomial equations.  For $t\ge0$, the norm bounds can be written as
\begin{align*}
 \opnorm{X}\le t
 &\quad\Longleftrightarrow\quad t^2I-X^*X\succeq0,\\
 \opnorm{X}\ge t
 &\quad\Longleftrightarrow\quad
 \exists v\;\bigl(\opnorm{v}=1\ \text{ and }\
                       \opnorm{Xv}^2\ge t^2\bigr).
\end{align*}
Positive semidefiniteness is determined by the principal minors.  After
splitting real and imaginary parts, all these conditions are polynomial
equalities and inequalities over $\mathbb R$.  So the existence of the matrices is a sentence in the
first-order theory of real closed fields, which is decidable by Tarski's theorem~\cite{Tarski}.  Let
$C(P,n,c)$ mean that $c$ encodes a valid certificate.  Then $C$ is
decidable.

Suppose first that $G_P$ is MF.  Put every trivial word in $W_n$ into $T$
and every remaining word into $S$.  Each word in $T$ has a normal-closure
expression of the required form.  Applying Lemma~\ref{lem:local-models}
to the finite family $S$, with relator tolerance $2^{-n}$, supplies the
matrices required for a valid certificate.

Conversely, suppose that a valid certificate exists for every $n$.  Let
$w_1,\ldots,w_q$ represent nonidentity elements of $G_P$, and let
$\varepsilon>0$.  Choose $n$ larger than the lengths of these words and
such that $2^{-n}\le\varepsilon$.  None of the $w_j$ can lie in $T$, since
membership in $T$ includes a proof that the word is trivial.  Thus every
$w_j$ lies in $S$, and the matrices in the certificate satisfy the
relator and separation bounds in Lemma~\ref{lem:local-models}, with
separation $1/4$ in place of $1/2$.  Korchagin's amplification argument
allows any fixed positive separation constant, so $G_P$ is MF.

So the displayed equivalence holds,
$\mathrm{MF}_{\mathrm{fp}}\in\Pi^0_2$, and
$\mathrm{NONMF}_{\mathrm{fp}}\in\Sigma^0_2$.
\end{proof}

\section{HNN extensions with a corona conjugator}\label{sec:hnn}

In~\cite{NonMF} we prove non-approximability using one-sided
conjugation in a norm matrix corona.  If the edge isomorphism
of an HNN extension is implemented by a
unitary of a norm matrix corona, then the HNN extension is MF\@.  For
$\theta=\mathrm{id}$ this is close to Shulman's theorem on central HNN
extensions~\cite{Shulman}; the version below assumes a tracial state
vanishing off the identity and allows a nontrivial edge isomorphism.  The
corona unitary implements the edge isomorphism, while the trace and the left
regular representation prove that the group embeds in the resulting
algebra.

A countable group $G$ \emph{admits a tracial MF realization}
$(A,\rho,\tau)$ if $A$ is a separable unital MF $C^*$-algebra,
$\rho\colon G\to\U(A)$ is a homomorphism, and $\tau$ is a tracial state on
$A$ with $\tau(\rho(g))=0$ for every $g\ne1$.  Since $\tau(1)=1$, such a
group embeds in $\U(A)$, so it is MF\@.  Only the vanishing of $\tau$
off the identity is used below; the constructions provide tracial
states.
\leanverified{Manuscript/OneSidedMFRadical/HNNCoronaConjugatorSentenceAudit}{GroupApproximation.Manuscript.OneSidedMFRadical.HNNCoronaConjugatorSentenceAudit.RegularRealizationData.rho_injective}
\leanverified{Manuscript/OneSidedMFRadical/HNNCoronaConjugatorSentenceAudit}{GroupApproximation.Manuscript.OneSidedMFRadical.HNNCoronaConjugatorSentenceAudit.RegularRealizationData.isOperatorMF}
To embed an arbitrary norm matrix corona $\mathcal Q_{\mathbf d}$ in
$\prod_kM_k(\C)/\bigoplus_kM_k(\C)$, choose a strictly increasing sequence
$(m_n)$ with $m_n\ge d_n$, embed $M_{d_n}(\C)$ as a corner of
$M_{m_n}(\C)$, and put zero in the unused coordinates.  The resulting map
on reduced products is an isometric $*$-homomorphism,
so the MF algebras of Section~\ref{sec:prelim} are exactly the
separable $C^*$-algebras that embed in the latter corona, as
in~\cite{BK,Shulman}.

\begin{lemma}[residually finite groups]\label{lem:rf-regular}
Every countable residually finite group admits a tracial MF realization.
\leanverified{Manuscript/OneSidedMFRadical/HNNCoronaConjugatorSentenceAudit}{GroupApproximation.Manuscript.OneSidedMFRadical.HNNCoronaConjugatorSentenceAudit.residuallyFinite_isRegularlyRealized}
\end{lemma}

\begin{proof}
Choose finite-index normal subgroups $N_1\ge N_2\ge\cdots$ with trivial
intersection, and let $\lambda_k$ be the left regular representation of
$G/N_k$, of dimension $d_k$.  In the corona $\mathcal Q_{\mathbf d}$ the
map $\rho(g)=[(\lambda_k(g))_k]$ is a homomorphism, and
$\tau([(x_k)_k])=\lim_\omega\operatorname{tr}_{d_k}(x_k)$, for a free
ultrafilter $\omega$, is a tracial state with $\tau(\rho(g))=0$ for
$g\ne1$, since $g\notin N_k$ for all large $k$.  The $C^*$-algebra
generated by $\rho(G)$ and the unit of the corona is separable and unital,
and the triple with $\rho$ and $\tau$ is the required realization.
\end{proof}

Let $G$ be a countable group, let $S\le G$ be a subgroup, and let
$\theta\colon S\to G$ be an injective homomorphism.  The HNN extension
\begin{equation}\label{eq:hnn}
 R=\langle G,t\mid tst^{-1}=\theta(s)\ (s\in S)\rangle
\end{equation}
contains $G$ by Britton's lemma~\cite{Britton}.
\leanverified{Manuscript/OneSidedMFRadical/HNNCoronaConjugatorSentenceAudit}{GroupApproximation.Manuscript.OneSidedMFRadical.HNNCoronaConjugatorSentenceAudit.hnnExtension_base_injective}

\begin{theorem}[HNN permanence with a corona conjugator]\label{thm:hnn-permanence}
In the situation of~\eqref{eq:hnn}, suppose that $G$ admits a tracial
MF realization $(A,\rho,\tau)$, and that for some norm matrix corona
$\mathcal Q$ there are an injective $*$-homomorphism
$\iota\colon A\to\mathcal Q$ and a unitary $W\in\mathcal Q$ with
\[
 W\,\iota\rho(s)\,W^*=\iota\rho(\theta(s))
 \qquad(s\in S).
\]
Then $R$ admits a tracial MF realization.  In particular $R$ is MF\@.
\end{theorem}

\begin{proof}
Write $D=C^*(\iota\rho(G))\subseteq\mathcal Q$, a separable
$C^*$-algebra with unit $p=\iota(1)$, and let
$B_0=C^*(\iota\rho(S))$ and $B_1=C^*(\iota\rho(\theta(S)))$, both
containing $p$.  Conjugation by $W$ restricts to a $*$-isomorphism
$\Theta\colon B_0\to B_1$ carrying $\iota\rho(s)$ to $\iota\rho(\theta(s))$.
Let $U$ be the quotient of the full unital free product $D*C(\mathbb T)$ by
the closed ideal generated by the elements $ubu^*-\Theta(b)$, $b\in B_0$,
where $u$ is the canonical unitary generator of $C(\mathbb T)$.  This is
the universal $C^*$-HNN extension of $D$ associated to $\Theta$; we use
$d$ and $u$ for the canonical images of $d\in D$ and the stable unitary.
\leanverified{Manuscript/OneSidedMFRadical/HNNCoronaConjugatorSentenceAudit}{GroupApproximation.Manuscript.OneSidedMFRadical.HNNCoronaConjugatorSentenceAudit.universalCStarHNN_mappingProperty}
Let
$C=B_0\oplus B_1$, $A_1=M_2(D)$, and $A_2=M_2(B_0)$, with the unital
inclusions
\[
 \iota_A(c_0,c_1)=\begin{pmatrix}c_0&0\\0&c_1\end{pmatrix}\in A_1,
 \qquad
 \iota_B(c_0,c_1)=\begin{pmatrix}c_0&0\\0&\Theta^{-1}(c_1)\end{pmatrix}\in A_2,
\]
and let $P=A_1*_CA_2$ be the full amalgamated free product, in which
$A_1$ and $A_2$ embed~\cite[Section~2.1]{Shulman}.  Ueda proved that the
universal HNN extension is a full corner of exactly this amalgamated free
product~\cite[Proposition~2.4]{Ueda}: for the projection
$e=\iota_A(1_D,0)=\iota_B(1_D,0)$ there is an isomorphism of $U$ onto
$ePe$ carrying $d\in D$ to $\operatorname{diag}(d,0)$ and the unitary
$u$ to the partial isometry $e_{12}f_{21}$ built from the matrix units
$e_{ij}$ of $A_1$ and $f_{ij}$ of $A_2$.  In particular $U$ embeds
in $P$.

Shulman's criterion reduces the MF property of $P$ to compatible
embeddings of $A_1$ and $A_2$ in one norm matrix
corona~\cite[Theorem~16]{Shulman}.  Identify $M_2(\mathcal Q)$ with the
norm matrix corona for the doubled dimension sequence.  Let $\varphi_A$
be the entrywise inclusion of $M_2(D)$, and let $\varphi_B$ be the
entrywise inclusion of $M_2(B_0)$ followed by conjugation by
$\operatorname{diag}(1,W)$.  Both maps are injective, and for
$(c_0,c_1)\in C$,
\[
 \varphi_B(\iota_B(c_0,c_1))
 =\begin{pmatrix}c_0&0\\0&W\Theta^{-1}(c_1)W^*\end{pmatrix}
 =\begin{pmatrix}c_0&0\\0&c_1\end{pmatrix}
 =\varphi_A(\iota_A(c_0,c_1)),
\]
because $\Theta$ is conjugation by $W$.  So $P$ is MF, as is its
$C^*$-subalgebra $ePe\cong U$.

The functional $\tau\circ\iota^{-1}$ is a tracial state on
$\iota(A)\supseteq D$ that vanishes at $\iota\rho(g)$ for $g\ne1$.  In
its GNS representation, the orbit of the cyclic vector under
$\iota\rho(G)$ is orthonormal and carries the left regular
representation of $G$.
\leanverified{Manuscript/OneSidedMFRadical/HNNCoronaConjugatorSentenceAudit}{GroupApproximation.Manuscript.OneSidedMFRadical.HNNCoronaConjugatorSentenceAudit.regularCharacterGNS_orbit_orthonormal}
\leanverified{Manuscript/OneSidedMFRadical/HNNCoronaConjugatorSentenceAudit}{GroupApproximation.Manuscript.OneSidedMFRadical.HNNCoronaConjugatorSentenceAudit.regularCharacterGNS_left_translation}
Since $\ell^2(R)$ is a direct sum of copies of $\ell^2(G)$ over the
right cosets of $G$, the assignment $\lambda_G(g)\mapsto\lambda_R(g)$
extends to $C^*_{\mathrm r}(G)\to C^*_{\mathrm r}(R)$; composing gives
$\sigma_0\colon D\to C^*_{\mathrm r}(R)$ with
$\sigma_0(\iota\rho(g))=\lambda_R(g)$.  The pair
$(\sigma_0,\lambda_R(t))$ is covariant for $\Theta$, because
$tst^{-1}=\theta(s)$ in $R$, so the universal property of $U$ gives
$\sigma\colon U\to C^*_{\mathrm r}(R)$ with
$\sigma(\iota\rho(g))=\lambda_R(g)$ and $\sigma(u)=\lambda_R(t)$.

Define $j\colon R\to\U(U)$ by $j|_G=\iota\rho$ and $j(t)=u$.  The
relations of~\eqref{eq:hnn} hold in $U$, so $j$ is a homomorphism, and
$\sigma\circ j=\lambda_R$ is injective, so $j$ is injective.
\leanverified{Manuscript/OneSidedMFRadical/HNNCoronaConjugatorSentenceAudit}{GroupApproximation.Manuscript.OneSidedMFRadical.HNNCoronaConjugatorSentenceAudit.universalCStarHNN_groupLift_injective}
Pulling the canonical trace of $C^*_{\mathrm r}(R)$ back along $\sigma$
makes $(C^*(j(R)),\,j,\,\tau_R\circ\sigma)$ a tracial MF realization
of $R$.
\end{proof}

Injectivity of $j$ follows from $\sigma\circ j=\lambda_R$, not from the
covariant pair $(\iota\rho,W)$, which need not be faithful; in the
applications below it is not.

\begin{corollary}[central HNN extensions]\label{cor:central-hnn}
If $G$ admits a tracial MF realization and $S\le G$ is any subgroup, then
$\langle G,t\mid[t,s]=1\ (s\in S)\rangle$ admits a tracial MF realization.
\end{corollary}

\begin{proof}
Apply Theorem~\ref{thm:hnn-permanence} with $\theta$ the inclusion of
$S$, with $\iota$ any injective $*$-homomorphism of $A$ into a norm matrix
corona, and with $W=1$.
\end{proof}

The positive branch of the reduction below produces tracial MF
realizations inside norm reduced products.  For $C^*$-algebras $B_n$, $n\ge1$, write
$\prod_nB_n/\bigoplus_nB_n$ for the quotient of the bounded sequences by
the sequences with $\opnorm{x_n}\to0$; its norm is $\limsup_n\opnorm{x_n}$.

\begin{lemma}[reduced products of MF algebras]\label{lem:reduced-products}
If every $B_n$ is MF, then every separable $C^*$-subalgebra of
$\prod_nB_n/\bigoplus_nB_n$ is MF\@.  Moreover $M_k(B)$ is MF whenever
$B$ is.
\leanverified{Analysis/MFAlgebraMatrixAmplification}{GroupApproximation.MFAlgebraMatrixAmplification.isMFAlgebra_cstarMatrix}
\end{lemma}

The first assertion follows by diagonalization from the matrix-corona
characterization of MF algebras (Blackadar--Kirchberg~\cite{BK}; see
also~\cite[Chapter~11]{BrownOzawa}), and Korchagin uses it in this
form~\cite[Proposition~6]{Korchagin}.  The second holds because $M_k$ of
a norm matrix corona is the norm matrix corona of the dimension sequence
multiplied by $k$.

\begin{lemma}[tensor synchronization]\label{lem:tensor-sync}
Let $\Gamma$ be a countable group with a tracial MF realization
$(A_1,\rho_1,\tau_1)$, let $Q$ be a countable group with homomorphisms
$\beta_n\colon Q\to B_n$ to finite groups such that every $q\ne1$
satisfies $\beta_n(q)\ne1$ for all large $n$, let $S\le\Gamma$, and let
$\tau\colon S\to Q$ be a homomorphism.  Suppose that for every $n$ there
is a homomorphism $\lambda_n\colon\Gamma\to G_n$ to a finite group with
\[
 \ker(\lambda_n|_S)\le\ker(\beta_n\circ\tau).
\]
Then there are a separable unital MF $C^*$-algebra $A$, a homomorphism
$V\colon\Gamma\times Q\to\U(A)$, a tracial state $T$ on $A$ with
$T(V(g,q))=0$ for $(g,q)\ne(1,1)$, and a unitary $W\in A$ with
$WV(s,1)W^*=V(s,\tau(s))$ for all $s\in S$.  Consequently
\[
 \bigl\langle\Gamma\times Q,\ u\ \bigm|\ u(s,1)u^{-1}=(s,\tau(s))\ (s\in S)\bigr\rangle
\]
admits a tracial MF realization, and in particular is MF\@.
\end{lemma}

\begin{proof}
Let $E_n$ be the image of $(\lambda_n,\beta_n)\colon\Gamma\times Q\to G_n\times B_n$,
a finite group of order $k_n$, and let $L_n$ be its left regular
representation on $\ell^2(E_n)$.  Put $B'_n=A_1\otimes M_{k_n}(\C)$,
which is MF by Lemma~\ref{lem:reduced-products}, let
$E=\prod_nB'_n/\bigoplus_nB'_n$, and define
\[
 V(g,q)=\bigl[\bigl(\rho_1(g)\otimes L_n(\lambda_n(g),\beta_n(q))\bigr)_n\bigr],
 \qquad
 T\bigl([(x_n)_n]\bigr)=\lim_\omega(\tau_1\otimes\operatorname{tr}_{k_n})(x_n),
\]
for a free ultrafilter $\omega$.  Then $V$ is a homomorphism and $T$
is a tracial state on $E$.  If $(g,q)\ne(1,1)$, then either $g\ne1$ and
$\tau_1(\rho_1(g))=0$, or $g=1$ and $\beta_n(q)\ne1$ for all large $n$
while the regular character of a finite group vanishes off the identity;
either way $T(V(g,q))=0$.

The two homomorphisms $s\mapsto(\lambda_n(s),1)$ and
$s\mapsto(\lambda_n(s),\beta_n\tau(s))$ from $S$ to $E_n$ have the same
kernel, namely $\ker(\lambda_n|_S)$, by the hypothesis.  So their images are
isomorphic subgroups of $E_n$ of the same order, isomorphic via
$(\lambda_n(s),1)\mapsto(\lambda_n(s),\beta_n\tau(s))$.  The restriction
of $L_n$ to a subgroup of $E_n$ is unitarily equivalent to the direct sum
of as many copies of the left regular representation of that subgroup as
it has right cosets, so there is a unitary $W_n\in M_{k_n}(\C)$ with
$W_nL_n(\lambda_n(s),1)W_n^*=L_n(\lambda_n(s),\beta_n\tau(s))$ for all
$s\in S$.  Put $W=[(1\otimes W_n)_n]\in E$; then $WV(s,1)W^*=V(s,\tau(s))$
for $s\in S$.  Let $A$ be the $C^*$-algebra generated by $V(\Gamma\times Q)$ and $W$;
it is separable, unital, and MF by Lemma~\ref{lem:reduced-products}, and
$(A,V,T|_A)$ is a tracial MF realization of $\Gamma\times Q$.  For the
last assertion, embed $A$ in a norm matrix corona by some injective
$\iota$, extend $\iota(W)$ to the unitary $\iota(W)+(1-\iota(1))$ of the
full corona, and apply Theorem~\ref{thm:hnn-permanence}.
\end{proof}

\section{The complexity of recognizing MF groups}\label{sec:compiler}

Write $W_e$ for the domain of the $e$-th partial computable function and
\[
 \mathrm{INF}=\{e:W_e\text{ is infinite}\},
 \qquad
 \mathrm{FIN}=\{e:W_e\text{ is finite}\}.
\]
The set $\mathrm{INF}$ is $\Pi^0_2$-complete and $\mathrm{FIN}$ is
$\Sigma^0_2$-complete~\cite[Chapter~IV]{Soare}.  The construction
combines the non-MF group $H$ of~\cite{NonMF} with Theorem~\ref{thm:hnn-permanence}: it
produces a computable map
$e\mapsto\widehat R_e$ into finite presentations such that the
presented group contains $H$ when $e\in\mathrm{FIN}$ and satisfies the
hypotheses of Theorem~\ref{thm:hnn-permanence} when $e\in\mathrm{INF}$,
so that it is
MF if and only if $e\in\mathrm{INF}$.
For $e\in\mathrm{FIN}$, the group $\widehat R_e$ is not MF; for
$e\in\mathrm{INF}$, Lemma~\ref{lem:positive-branch} proves that it is
MF.
Table~\ref{tab:objects} lists the objects of the construction.

\begin{table}[htb]
\centering\small
\begin{tabular}{@{}l@{\qquad}p{0.66\linewidth}@{}}
$H$ & the non-MF group $\EL_{12}(L_{\F_2}(1,2))$ of~\cite{NonMF} \\[1pt]
$C_e$ & recursively presented; trivial for $e\in\mathrm{INF}$, contains
  $H$ for $e\in\mathrm{FIN}$ \\[1pt]
$Q_e=F/N_e$ & three-generated recursively presented group containing
  $C_e$; here $F=F(x,y,t)$ \\[1pt]
$Q_+=F/N_+$ & the value of $Q_e$ for trivial $C_e$; embeds in
  $F(x_1,y)\times F(x_2,t)$ \\[1pt]
$\Lambda_e$; $M_e$ & finitely presented Higman host of $Q_e$; the
  Mihailova subgroup detecting equality in $\Lambda_e$ \\[1pt]
$K_e$; $L_e$ & a product of six free groups; its finitely generated
  subgroup with $i(F)\cap L_e=i(N_e)$ \\[1pt]
$\Gamma_e$; $S_e$ & the central HNN extension
  $\langle K_e,v\mid[v,\ell]=1\ (\ell\in L_e)\rangle$;
  $S_e=\langle i(F),v\,i(F)\,v^{-1}\rangle\cong F*_{N_e}F$ \\[1pt]
$\alpha_e\colon S_e\to Q_e$ & $q_e$ on $i(F)$, trivial on
  $v\,i(F)\,v^{-1}$ \\[1pt]
$\widehat R_e$ & the finite presentation of
  $\langle\Gamma_e\times Q_e,\ u\mid
   u(s,1)u^{-1}=(s,\alpha_e(s))\ (s\in S_e)\rangle$
\end{tabular}
\caption{The objects of the construction, in order of appearance.}
\label{tab:objects}
\end{table}

\subsection*{A recursive presentation of \texorpdfstring{$H$}{H}}

\begin{lemma}[a recursive presentation of $H$]\label{lem:seed}
The group $H=\EL_{12}(L_{\F_2}(1,2))$ is finitely generated and has decidable word
problem.  In particular it has a recursive presentation on finitely many
generators.
\end{lemma}

\begin{proof}
Let $Y$ be the set of elements $e_{ij}(a)$ with $i\ne j$ and
$a\in\{1,s_0,s_1,t_0,t_1\}$.  Since
$e_{ij}(a+b)=e_{ij}(a)e_{ij}(b)$ and
$[e_{ik}(a),e_{kj}(b)]=e_{ij}(ab)$ for distinct $i,j,k$, induction on
monomial length shows that $H=\langle Y\rangle$.  The monomials in
$s_0,s_1,t_0,t_1$ that contain no subword $t_is_j$ and no subword
$s_1t_1$ form the standard reduced basis of $L_{\F_2}(1,2)$ by the basis
theorem for Leavitt path algebras, with $s_1$ chosen as the special edge
\cite[Theorem~3.7]{LopatkinNam}.  The reductions
$t_is_j\to\delta_{ij}$ and $s_1t_1\to1+s_0t_0$ compute products in this
basis, so equality in $L_{\F_2}(1,2)$ is decidable.  Thus a word in $Y^{\pm1}$
evaluates to an explicit $12\times12$ matrix over
$L_{\F_2}(1,2)$, and it represents $1$ in $H$ if and only if
that matrix is the identity.  The set of words representing $1$ is
decidable, so recursively enumerable, and $\langle Y\mid\text{that set}\rangle$
is a recursive presentation of $H$.
\end{proof}

\subsection*{The group \texorpdfstring{$C_e$}{C-e}}

\begin{lemma}[the group $C_e$]\label{lem:switch}
There is an algorithm that computes from $e$ a recursive presentation, on
a computable set of generators, of a group $C_e$ such that $C_e$ is
trivial if $e\in\mathrm{INF}$ and $H$ embeds in $C_e$ if
$e\in\mathrm{FIN}$.
\end{lemma}

\begin{proof}
Apply the construction of Bilanovic--Chubb--Roven~\cite[Theorem~3.1]{BCR}
to the trivial group and $H$.
Let $\langle Y\mid T\rangle$ be the recursive presentation of $H$ from
Lemma~\ref{lem:seed}.  Take generators $y_{i,\ell}$ for $i\ge1$ and
$\ell\in Y$, one copy $Y_i$ of $Y$ for each $i$, and as relators all
words of $T$ rewritten in each copy $Y_i$, together with every generator
$y_{i,\ell}$ for which $i\le\lvert W_e\rvert$.  The last family is
recursively enumerable uniformly in $e$: enumerate $W_e$, and whenever its
$i$-th element appears, enumerate all of $Y_i$.  The group presented is
the free product of the copies $H_i\cong H$, $i\ge1$, with the first
$\lvert W_e\rvert$ copies killed.  If $W_e$ is infinite every copy is
killed and $C_e$ is trivial; if $W_e$ is finite, $C_e$ is the free product
of the remaining copies and contains $H$.
\end{proof}

\subsection*{A three-generator embedding}

Let $F=F(x,y,t)$ be the free group on $x,y,t$, and let
$y_i=x^iyx^{-i}$ for $i\in\mathbb Z$.  For a countable group $C$ with a
generating sequence $(c_i)_{i\ge1}$, put $c_i=1$ for $i\le0$ and
\begin{equation}\label{eq:bridge}
 B(C)=\bigl\langle C*F(x,y),\ t\ \bigm|\ ty_it^{-1}=c_iy_i\ (i\in\mathbb Z)\bigr\rangle.
\end{equation}
Let $Q_+=B(1)=\langle x,y,t\mid[t,y_i]=1\ (i\in\mathbb Z)\rangle$, let
$q_+\colon F\to Q_+$ be the quotient map, and let $N_+=\ker q_+$.

\begin{lemma}[the three-generator embedding]\label{lem:bridge}
\begin{enumerate}
\item $B(C)$ is an HNN extension of $C*F(x,y)$, so $C$ embeds in $B(C)$;
it is generated by $x,y,t$; and a recursive presentation of $B(C)$ on
$x,y,t$ is computable from a recursive presentation of $C$.  We write
$B(C)=F/N_C$ and $q_C\colon F\to B(C)$ for the quotient map.
\item Killing $C$ induces a surjection $B(C)\to Q_+$; so
$N_C\le N_+$, with equality when $C$ is trivial.
\item The assignment $x\mapsto(x_1,x_2)$, $y\mapsto(y,1)$, $t\mapsto(1,t)$
defines an injective homomorphism $j\colon Q_+\to P=F(x_1,y)\times F(x_2,t)$.
In particular $Q_+$ is residually finite.
\end{enumerate}
\end{lemma}

\begin{proof}
(1) The elements $y_i$, $i\in\mathbb Z$, freely generate the normal
closure of $y$ in $F(x,y)$~\cite[Chapter~I]{LyndonSchupp}, and the
retraction of $C*F(x,y)$ onto $F(x,y)$ carries a nontrivial reduced word
in the $c_iy_i$ to a nontrivial word in the $y_i$; so $y_i\mapsto c_iy_i$
is an isomorphism between two free subgroups of the base,
and~\eqref{eq:bridge} is an HNN extension.  Britton's lemma embeds the
base, and with it $C$.  The relation $c_i=ty_it^{-1}y_i^{-1}$ lets
$x,y,t$ generate $B(C)$, and substituting it into the relators produces
a recursive presentation on $x,y,t$, uniformly in the presentation
of $C$.

(2) The relators of~\eqref{eq:bridge} map to the relators $[t,y_i]$ of
$Q_+$ when every $c_i$ is sent to $1$, and the map is the identity on
$x,y,t$.

(3) The relators $[t,y_i]$ die in $P$, so $j$ is a homomorphism.  Let
$N_0$ be the normal closure of $\{y,t\}$ in $Q_+$; then
$Q_+=N_0\rtimes\langle x\rangle$, and rewriting along the transversal
$\{x^n\}$ by the Reidemeister--Schreier
method~\cite[Chapter~II]{LyndonSchupp} presents $N_0$ as
$F(y_n:n\in\mathbb Z)\times F(t_m:m\in\mathbb Z)$, where
$y_n=x^nyx^{-n}$ and $t_m=x^mtx^{-m}$.  The map $j$ carries these two
free factors isomorphically onto the normal closures of $y$ in $F(x_1,y)$
and of $t$ in $F(x_2,t)$, and the exponent sum of $x_1$ reads off the
$\langle x\rangle$-coordinate; so $j$ is injective.  Free groups are
residually finite, and residual finiteness passes to direct products and
subgroups, so $Q_+$ is residually finite.
\end{proof}

For $e\in\mathbb N$ let $Q_e=B(C_e)=F/N_e$ with $C_e$ the group of
Lemma~\ref{lem:switch}, and write $q_e\colon F\to Q_e$.  By
Lemmas~\ref{lem:switch} and~\ref{lem:bridge}, a recursive presentation of
$Q_e$ on $x,y,t$ is computable from $e$; $N_e\le N_+$; if $e\in\mathrm{INF}$
then $Q_e=Q_+$, $N_e=N_+$ and $q_e=q_+$; and if $e\in\mathrm{FIN}$ then
$H\le C_e\le Q_e$.

\subsection*{Detecting \texorpdfstring{$N_e$}{N-e} inside a finitely presented group}

\begin{lemma}[Higman embedding and the Mihailova subgroup]\label{lem:mikhailova}
There is an algorithm that computes from $e$ a finite presentation
$\langle X_e\mid R_e\rangle$ of a group $\Lambda_e$, words
$w_x,w_y,w_t\in F(X_e)$, and a finite generating set of a subgroup
$M_e\le F(X_e)\times F(X_e)$, such that $x\mapsto w_x$, $y\mapsto w_y$,
$t\mapsto w_t$ induces an embedding $Q_e\to\Lambda_e$, and
$M_e=\{(u,v):u=v\text{ in }\Lambda_e\}$.  So, with
\[
 K^0_e=F\times F(X_e)\times F(X_e),
 \qquad
 i_e(f)=(f,w_f,1),
 \qquad
 L^0_e=F\times M_e,
\]
where $w_f$ is the image of $f$ under $x\mapsto w_x$, $y\mapsto w_y$,
$t\mapsto w_t$, the map $i_e\colon F\to K^0_e$ is injective and
$i_e(f)\in L^0_e$ if and only if $f\in N_e$.
\end{lemma}

\begin{proof}
Higman's embedding theorem~\cite{Higman} embeds every finitely generated
recursively presented group in a finitely presented group, and its proof
is effective; an explicit algorithm producing the finite presentation and
the embedding words from a recursive presentation is given by
Mikaelian~\cite{Mikaelian}.  Apply it to the presentation of $Q_e$.
Let $M_e$ be the subgroup of $F(X_e)\times F(X_e)$ generated by the
pairs $(z,z)$, $z\in X_e$, and $(r,1)$, $r\in R_e$; then
$M_e=\{(u,v):u=v\text{ in }\Lambda_e\}$ by Mihailova's
lemma~\cite{Mihailova}.  The first coordinate makes $i_e$
\leanverified{Manuscript/OneSidedMFRadical/ComputabilityConstructionClosure}{GroupApproximation.Manuscript.OneSidedMFRadical.ComputabilityConstruction.manuscriptPrintedMihailovaFiberProductPackage}
injective, and $(f,w_f,1)\in F\times M_e$ if and only if $w_f=1$ in
$\Lambda_e$, if and only if $q_e(f)=1$ because the embedding is
injective, if and only if $f\in N_e$.
\end{proof}

\subsection*{The two HNN extensions}

Let $P=F(x_1,y)\times F(x_2,t)$ and $j\colon Q_+\to P$ be as in
Lemma~\ref{lem:bridge}, and put
\[
 K^{\mathrm g}=F\times P,
 \qquad
 L^{\mathrm g}=\{(f,\,jq_+(f)):f\in F\},
\]
and define
\[
 K_e=K^0_e\times K^{\mathrm g},
 \qquad
 L_e=L^0_e\times L^{\mathrm g},
 \qquad
 i\colon F\to K_e,\quad i(f)=\bigl(i_e(f),(f,1)\bigr).
\]
Let
\begin{equation}\label{eq:central-rope}
 \Gamma_e=\bigl\langle K_e,\ v\ \bigm|\ [v,\ell]=1\ (\ell\in L_e)\bigr\rangle,
 \qquad
 S_e=\bigl\langle i(F),\ v\,i(F)\,v^{-1}\bigr\rangle\le\Gamma_e.
\end{equation}

\begin{lemma}[the first HNN extension]\label{lem:central-rope}
\begin{enumerate}
\item $K_e$ is a direct product of free groups of finite rank, so
finitely presented and residually finite; $L_e$ is finitely generated;
$i$ is injective; and $i(F)\cap L_e=i(N_e)$.
\item $\Gamma_e$ is finitely presented, and a finite presentation is
computable from $e$.  The natural map from the amalgamated free product
$F*_{N_e}F$, in which $n\in N_e$ in the first copy is identified with
$n$ in the second copy, onto $S_e$, sending the first copy to $i(F)$ and
the second to $v\,i(F)\,v^{-1}$, is an isomorphism.
\item There is a homomorphism $\alpha_e\colon S_e\to Q_e$ with
$\alpha_e(i(f))=q_e(f)$ and $\alpha_e(v\,i(f)\,v^{-1})=1$ for $f\in F$.
\end{enumerate}
\end{lemma}

\begin{proof}
(1) $K_e$ is a product of six free groups of finite rank, so it is
finitely presented and residually finite, and $L_e$ is generated by the
basis of the first factor, the generators of $M_e$ from
Lemma~\ref{lem:mikhailova}, and the elements $(a,jq_+(a))$,
$a\in\{x,y,t\}$.  Injectivity of $i$ is read off the first coordinate.
If $i(f)\in L_e$, then the Mihailova coordinate forces $f\in N_e$ by
Lemma~\ref{lem:mikhailova} and the graph coordinate forces $f\in N_+$;
since $N_e\le N_+$, we get $i(F)\cap L_e=i(N_e)$.

(2) The presentation~\eqref{eq:central-rope} is finite once $[v,\ell]=1$ is
imposed only for the finitely many generators $\ell$ of $L_e$, and it is
computable from $e$ because $K_e$, $L_e$ and $i$ are.  It is an HNN
extension of $K_e$ with both edge groups equal to $L_e$ and the identity
as edge isomorphism.  The map from $F*_{N_e}F$ to $S_e$ is well defined
because $i(n)\in L_e$ commutes with $v$ for $n\in N_e$.  A reduced word of
the amalgamated free product alternates between elements of the two
copies of $F$ lying outside $N_e$; its image is a word
$i(f_0)\,v\,i(f_1)\,v^{-1}\,i(f_2)\cdots$ in which every letter between
consecutive occurrences of $v^{\pm1}$ is some $i(f_k)$ with
$f_k\notin N_e$, so $i(f_k)\notin L_e$ by (1).  Such a word contains no
pinch $v^{\pm1}\ell v^{\mp1}$ with $\ell\in L_e$, so it is nontrivial in
$\Gamma_e$ by Britton's lemma.

(3) Define $\alpha_e$ on $F*_{N_e}F$ by $q_e$ on the first copy and by the
trivial map on the second; both kill $N_e$, so the definition is
consistent on the amalgamated subgroup, and (2) identifies $F*_{N_e}F$
with $S_e$.
\end{proof}

The second HNN extension is
\begin{equation}\label{eq:twisted-rope}
 R_e=\bigl\langle\Gamma_e\times Q_e,\ u\ \bigm|\ u(s,1)u^{-1}=(s,\alpha_e(s))\ (s\in S_e)\bigr\rangle,
\end{equation}
whose edge maps $s\mapsto(s,1)$ and $s\mapsto(s,\alpha_e(s))$ are injective
homomorphisms of $S_e$ into $\Gamma_e\times Q_e$; so
$\Gamma_e\times Q_e$ embeds in $R_e$ by Britton's lemma.  The
presentation~\eqref{eq:twisted-rope} lists the recursively enumerable
relators of $Q_e$.  Its finite replacement is
\begin{equation}\label{eq:finite-rope}
\begin{aligned}
 \widehat R_e=\Bigl\langle\Gamma_e\times F,\ u\ \Bigm|\ 
 &u\bigl(i(a),1\bigr)u^{-1}=\bigl(i(a),a\bigr),\\
 &u\bigl(v\,i(a)\,v^{-1},1\bigr)u^{-1}=\bigl(v\,i(a)\,v^{-1},1\bigr)
 \quad(a\in\{x,y,t\})\Bigr\rangle,
\end{aligned}
\end{equation}
where $\Gamma_e\times F$ is equipped with the finite presentation of
Lemma~\ref{lem:central-rope}(2), the basis $x,y,t$ of $F$, and the
commutation relators between the two factors.

\begin{lemma}[the finite presentation]\label{lem:finite-rope}
The groups $\widehat R_e$ and $R_e$ are isomorphic by the map that is the
identity on $\Gamma_e$ and on $u$ and sends $F$ onto $Q_e$ by $q_e$.  A
code of the finite presentation~\eqref{eq:finite-rope} is computable
from $e$.
\end{lemma}

\begin{proof}
Both maps $f\mapsto u(i(f),1)u^{-1}$ and $f\mapsto(i(f),f)$ are
homomorphisms $F\to\widehat R_e$, and they agree on the basis, so
$u(i(f),1)u^{-1}=(i(f),f)$ for all $f\in F$; likewise
$u(v\,i(f)\,v^{-1},1)u^{-1}=(v\,i(f)\,v^{-1},1)$ for all $f$.  If
$n\in N_e$, then $i(n)\in L_e$ by Lemma~\ref{lem:central-rope}(1), so
$v\,i(n)\,v^{-1}=i(n)$ in $\Gamma_e$, and comparing the two relations
gives $(i(n),n)=(i(n),1)$, that is, $(1,n)=1$ in $\widehat R_e$.  So
the factor $F$ of $\widehat R_e$ factors through $Q_e=F/N_e$, and the
relations of~\eqref{eq:finite-rope} become $u(s,1)u^{-1}=(s,\alpha_e(s))$
for $s$ in the generating set $\{i(a),\,v\,i(a)\,v^{-1}:a\in\{x,y,t\}\}$
of $S_e$; both sides define homomorphisms on $S_e$, so the relations hold
for all $s\in S_e$.  This defines a homomorphism $R_e\to\widehat R_e$; conversely, the
relations of~\eqref{eq:finite-rope} hold in $R_e$ after $F\to Q_e$,
because $\alpha_e(i(a))=q_e(a)$ and $\alpha_e(v\,i(a)\,v^{-1})=1$.  The two
homomorphisms are mutually inverse on generators.  Computability of the
code follows from Lemmas~\ref{lem:mikhailova} and~\ref{lem:central-rope}.
\end{proof}

\begin{lemma}[the case $e\in\mathrm{FIN}$]\label{lem:negative-branch}
If $e\in\mathrm{FIN}$, then $H$ embeds in $\widehat R_e$, and
$\widehat R_e$ is not MF\@.
\end{lemma}

\begin{proof}
For $e\in\mathrm{FIN}$ we have $H\le C_e\le Q_e\le\Gamma_e\times Q_e\le R_e\cong\widehat R_e$
by Lemmas~\ref{lem:switch},~\ref{lem:bridge}(1), and~\ref{lem:finite-rope}, and
Britton's lemma for~\eqref{eq:twisted-rope}.  Restricting MF models to a
subgroup shows that subgroups of MF groups are MF, and $H$ is not MF
by~\cite[Theorem~B]{NonMF}.
\leanverified{Manuscript/OneSidedMFRadical/ComputabilityConstructionClosure}{GroupApproximation.Manuscript.OneSidedMFRadical.ComputabilityConstruction.manuscriptPrintedNonMFEmbeddingObstruction}
\end{proof}

\begin{lemma}[the case $e\in\mathrm{INF}$]\label{lem:positive-branch}
If $e\in\mathrm{INF}$, then $\widehat R_e$ admits a tracial MF
realization; in particular, it is MF\@.
\end{lemma}

\begin{proof}
Let $e\in\mathrm{INF}$, so that $Q_e=Q_+$, $N_e=N_+$, $q_e=q_+$, and
$\alpha_e(i(f))=q_+(f)$, $\alpha_e(v\,i(f)\,v^{-1})=1$.  By
Lemma~\ref{lem:finite-rope} it suffices to show that $R_e$ admits a
tracial MF realization, and by Lemma~\ref{lem:tensor-sync} applied to $\Gamma=\Gamma_e$,
$Q=Q_+$, $S=S_e$ and $\tau=\alpha_e$, it suffices to produce the data of
that lemma.

The group $K_e$ is residually finite by Lemma~\ref{lem:central-rope}(1),
so it admits a tracial MF realization by Lemma~\ref{lem:rf-regular};
and $\Gamma_e$, the central HNN extension of $K_e$ over $L_e$, admits
one by Corollary~\ref{cor:central-hnn}.

The group $P=F(x_1,y)\times F(x_2,t)$ is residually finite; choose
finite-index normal subgroups $P_1\ge P_2\ge\cdots$ of $P$ with trivial
intersection, let $r_n\colon P\to C_n=P/P_n$, and put
$\beta_n=r_n\circ j\colon Q_+\to C_n$ with $j$ the embedding of
Lemma~\ref{lem:bridge}(3).  Since $j$ is injective and the $P_n$ decrease
to the trivial group, every $h\ne1$ has $\beta_n(h)\ne1$ for all large $n$.

Let $G_n=(C_n\times C_n)\rtimes\langle\sigma\rangle$, with $\sigma$ of
order two exchanging the coordinates.  Define $\lambda_n$ on
$K_e=K^0_e\times(F\times P)$ by killing $K^0_e$ and sending $(f,p)$ to
$(r_njq_+(f),\,r_n(p))$, and put $\lambda_n(v)=\sigma$.  The factor
$L^0_e$ of $L_e$ is killed, and $L^{\mathrm g}$ lands on the diagonal,
which commutes with $\sigma$; so $\lambda_n$ is a homomorphism on
$\Gamma_e$.  On $S_e\cong F*_{N_+}F$ it sends the first copy of $F$ into
the first coordinate and its $v$-conjugate into the second:
$\lambda_n|_{S_e}=(r_nj\pi_0,\,r_nj\pi_1)$, where $\pi_0=\alpha_e$ and
$\pi_1$ is $q_+$ on the second copy and trivial on the first, as in
Lemma~\ref{lem:central-rope}(3).  So
\[
 \ker(\lambda_n|_{S_e})=\ker(r_nj\pi_0)\cap\ker(r_nj\pi_1)
 \le\ker(\beta_n\circ\alpha_e),
\]
the hypothesis of Lemma~\ref{lem:tensor-sync}.
\end{proof}

\begin{proof}[Proof of Theorem~\ref{thm:recognition}]
By Lemma~\ref{lem:finite-rope} the map $e\mapsto\widehat R_e$ is
computable, by Lemma~\ref{lem:positive-branch} the presented group is MF
for $e\in\mathrm{INF}$, and by Lemma~\ref{lem:negative-branch} it is not MF
for $e\in\mathrm{FIN}$.  So $\mathrm{INF}$ reduces to
$\mathrm{MF}_{\mathrm{fp}}$ and $\mathrm{FIN}$ reduces to
$\mathrm{NONMF}_{\mathrm{fp}}$ under computable many-one reductions.  Since
$\mathrm{INF}$ is $\Pi^0_2$-complete and $\mathrm{FIN}$ is
$\Sigma^0_2$-complete, and since
$\mathrm{MF}_{\mathrm{fp}}\in\Pi^0_2$ and
$\mathrm{NONMF}_{\mathrm{fp}}\in\Sigma^0_2$ by
Proposition~\ref{prop:mf-upper-bound}, the two sets are complete at their
levels.  A $\Pi^0_2$-complete set is not $\Sigma^0_2$ and a
$\Sigma^0_2$-complete set is not $\Pi^0_2$, so neither set is recursively
enumerable, and in particular neither is decidable.
\end{proof}

\begin{remark}\label{rem:template}
The construction uses only these properties of $H$: it is finitely
generated, recursively presented, and not MF, while subgroups of MF
groups are MF\@.  So it applies to any group property that is
closed under subgroups, given a finitely generated recursively presented
group without the property and an analogue of
Lemma~\ref{lem:tensor-sync} for the twisted HNN extensions; the open
problems below record which input is missing for soficity and for
hyperlinearity.
\end{remark}

The homomorphism $V$ of Lemma~\ref{lem:tensor-sync} does not extend
faithfully to $R_e$ in the corona: the finite maps $\lambda_n$ are trivial
on the Mihailova factor, so $W$ commutes with $V(\gamma,1)$ for every
$\gamma\in K^0_e$, although $u$ does not commute with $(\gamma,1)$ in
$R_e$.  Injectivity is obtained in Theorem~\ref{thm:hnn-permanence} from
the universal algebra.

\subsection*{Open problems}

We do not know
whether the models of $\widehat R_e$ can be chosen to converge to
$C^*_{\mathrm r}(\widehat R_e)$ in operator norm, so that the reduced
$C^*$-algebra itself is MF; the argument produces only a tracial state on
an abstract MF algebra.  We do not know whether
$\widehat R_e$ is hyperlinear for $e\in\mathrm{INF}$; placing
hyperlinearity recognition at the same level would also require a
finitely presented non-hyperlinear group, and none is known.  We do not know whether an analogous family of finitely presented groups
exists for soficity; such a construction would require a
permutation-model replacement for Lemma~\ref{lem:tensor-sync}.  Finally,
we do not know the complexity of MF recognition under the strong
convergence convention of~\cite{GaoEtAl,Schafhauser}: the norms
$\opnorm{\lambda(w)}$ in the left regular representation are not in
general computable from a finite presentation, so
Proposition~\ref{prop:mf-upper-bound} does not apply.



\section*{Acknowledgments}

The author thanks Francesco Fournier-Facio for early advice and detailed
comments on earlier drafts, and Alon Dogon for organizing the ensuing
discussion and for comments on the surrounding literature.  The author also
thanks Tatiana Shulman for clarifying the terminology and historical context
of the MF question, and Caleb Eckhardt for discussions of MF representations
and matrix norms.

\section*{Use of AI and formal methods}

Claude Fable~5 and GPT-5.6 Sol assisted with mathematical exploration,
proof development, literature searches, Lean formalization, drafting, and
editing.  Lean~4 source is available at \coderepo.

\begin{thebibliography}{99}
\bibitem[NMF]{NonMF}
Sauers,
\emph{Non-MF Groups},
2026, Palomar registry entry
\url{https://palomar-registry.org/entry?id=PALOMAR-2026-08-24-000006&version=1}.

\bibitem[Ad]{Adian}
S.~I. Adian,
\emph{Algorithmic unsolvability of problems of recognition of certain properties of groups},
Dokl. Akad. Nauk SSSR \textbf{103} (1955), 533--535.

\bibitem[BCR]{BCR}
I.~Bilanovic, J.~Chubb, and S.~Roven,
\emph{Detecting properties from descriptions of groups},
Arch. Math. Logic \textbf{59} (2020), no.~3--4, 293--312,
\doi{10.1007/s00153-019-00690-x}.

\bibitem[BK]{BK}
B.~Blackadar and E.~Kirchberg,
\emph{Generalized inductive limits of finite-dimensional $C^*$-algebras},
Math. Ann. \textbf{307} (1997), no.~3, 343--380,
\doi{10.1007/s002080050039}.

\bibitem[Bri]{Britton}
J.~L. Britton,
\emph{The word problem},
Ann. of Math. (2) \textbf{77} (1963), 16--32.

\bibitem[BO]{BrownOzawa}
N.~P. Brown and N.~Ozawa,
\emph{$C^*$-algebras and finite-dimensional approximations},
Graduate Studies in Mathematics \textbf{88},
American Mathematical Society, Providence, RI, 2008.

\bibitem[CDE]{CDE}
J.~R. Carri\'on, M.~Dadarlat, and C.~Eckhardt,
\emph{On groups with quasidiagonal $C^*$-algebras},
J. Funct. Anal. \textbf{265} (2013), no.~1, 135--152,
\doi{10.1016/j.jfa.2013.04.004}.

\bibitem[GKM]{GaoEtAl}
D.~Gao, S.~Kunnawalkam Elayavalli, A.~Manzoor, and G.~Patchell,
\emph{A new source of purely finite matricial fields},
preprint, 2026,
\href{https://arxiv.org/abs/2603.24502}{arXiv:2603.24502}.

\bibitem[Hig]{Higman}
G.~Higman,
\emph{Subgroups of finitely presented groups},
Proc. Roy. Soc. London Ser. A \textbf{262} (1961), 455--475.

\bibitem[Kor]{Korchagin}
A.~I. Korchagin,
\emph{The MF-property for countable discrete groups},
Math. Notes \textbf{102} (2017), no.~2, 198--211,
\doi{10.1134/S0001434617070227},
\href{https://arxiv.org/abs/1704.06906}{arXiv:1704.06906}.

\bibitem[Lev]{Leavitt}
W.~G.~Leavitt,
\emph{The module type of a ring},
Trans. Amer. Math. Soc. \textbf{103} (1962), 113--130,
\doi{10.1090/S0002-9947-1962-0132764-X}.

\bibitem[Lem]{Lempp}
S.~Lempp,
\emph{The computational complexity of torsion-freeness of finitely presented groups},
Bull. Austral. Math. Soc. \textbf{56} (1997), 273--277.

\bibitem[LN]{LopatkinNam}
V.~Lopatkin and T.~G. Nam,
\emph{On the homological dimensions of Leavitt path algebras with coefficients in commutative rings},
J. Algebra \textbf{481} (2017), 273--292,
\doi{10.1016/j.jalgebra.2017.02.027}.

\bibitem[LS]{LyndonSchupp}
R.~C. Lyndon and P.~E. Schupp,
\emph{Combinatorial group theory},
Classics in Mathematics, Springer-Verlag, Berlin, 2001.

\bibitem[Mik]{Mikaelian}
V.~H. Mikaelian,
\emph{An explicit algorithm for the Higman embedding theorem},
preprint, 2025,
\href{https://arxiv.org/abs/2507.04347v8}{arXiv:2507.04347v8}.

\bibitem[Mih]{Mihailova}
K.~A. Mihailova,
\emph{The occurrence problem for direct products of groups},
Dokl. Akad. Nauk SSSR \textbf{119} (1958), 1103--1105.

\bibitem[Ra]{Rabin}
M.~O. Rabin,
\emph{Recursive unsolvability of group theoretic problems},
Ann. of Math. (2) \textbf{67} (1958), 172--194.

\bibitem[Sch]{Schafhauser}
C.~Schafhauser,
\emph{Finite dimensional approximations of certain amalgamated free
products of groups},
preprint, 2023,
\href{https://arxiv.org/abs/2306.02498}{arXiv:2306.02498}.

\bibitem[Sh]{Shulman}
T.~Shulman,
\emph{The MF property for amalgamated free products},
preprint, 2026,
\href{https://arxiv.org/abs/2603.13564v2}{arXiv:2603.13564v2}.

\bibitem[So]{Soare}
R.~I. Soare,
\emph{Recursively enumerable sets and degrees},
Perspectives in Mathematical Logic, Springer-Verlag, Berlin, 1987.

\bibitem[Ta]{Tarski}
A.~Tarski,
\emph{A decision method for elementary algebra and geometry},
2nd ed., University of California Press, Berkeley, 1951.

\bibitem[Ue]{Ueda}
Y.~Ueda,
\emph{Remarks on HNN extensions in operator algebras},
Illinois J. Math. \textbf{52} (2008), no.~3, 705--725,
\href{https://arxiv.org/abs/math/0601706}{arXiv:math/0601706}.

\end{thebibliography}
\end{document}
